\section{Three Middleware That Exist So a Request Will Not Run}
\label{sec:ch16-the-middleware-that-will-not-run}
\index{request pipeline!middleware that refuse|(}%

Chapter~\ref{ch:the-life-of-a-request} printed thirteen middleware and explained
twelve of them. The one it listed and left alone was
\code{SkipWarmupExecutionMiddleware} at position ten, which is a good name for
something whose purpose is not obvious.
Chapter~\ref{ch:identity-and-authorization} then added two more to three
services and located them by reading a startup log, which told it where they sat
and nothing about what they do. All three are in this section, and the last of
them turns out to have done something neither chapter noticed.

\subsection*{Position ten, and the request that is thrown away}

\index{warmup requests!the middleware at position ten}%
Here is the whole of \code{SkipWarmupExecutionMiddleware.cs}:

\begin{minted}[fontsize=\footnotesize,breaklines]{csharp}
public async ValueTask InvokeAsync(RequestContext context)
{
    if (context.IsWarmupRequest())
    {
        context.Result = new WarmupExecutionResult();
        return;
    }

    await next(context).ConfigureAwait(false);
}
\end{minted}

The interesting part is not the code, it is the address. Position ten sits after
the document cache, the parser, validation, the cost analyzer, the operation
cache and the compiler, and before variable coercion and execution. So a warmup
request is parsed, validated, costed and compiled, and then dropped on the
floor. It fills the caches and touches no resolver, no variable coercion and no
database. Everything expensive about a first request except the part that has
side effects.

\index{warmup requests!registering one}%
You get one by registering a warmup task.
\code{RequestExecutorManager.Warmup.cs} runs them when the executor is created:

\begin{minted}[fontsize=\footnotesize,breaklines]{csharp}
var warmupTasks = executor.Schema.Services
    .GetServices<IRequestExecutorWarmupTask>();

if (!isInitialCreation)
{
    warmupTasks = warmupTasks.Where(t => !t.ApplyOnlyOnStartup);
}

foreach (var warmupTask in warmupTasks)
{
    await warmupTask.WarmupAsync(executor, cancellationToken).ConfigureAwait(false);
}
\end{minted}

They run one after another rather than together, and \code{ApplyOnlyOnStartup}
distinguishes the first creation from a later one, which is what a schema reload
produces. Note where they are resolved from: the schema service provider, not
the application's. That is the same distinction that obliged
\code{MosaicSubgraphDefaults} to hand the logger factory across with
\code{AddApplicationService}, so that
chapter~\ref{ch:the-life-of-a-request}'s listener could find one, and it bites
in the same way.

\index{warmup requests!marked in process only}%
The marker, in \code{OperationRequestBuilderExtensions.cs}, is worth one
paragraph because it looks alarming and is not:

\begin{minted}[fontsize=\footnotesize,breaklines]{csharp}
/// Marks this request as a warmup request that will bypass security measures and skip execution.
public static OperationRequestBuilder MarkAsWarmupRequest(
    this OperationRequestBuilder builder)
    => builder.SetGlobalState(ExecutionContextData.IsWarmupRequest, true);
\end{minted}

It is set through global state on a request builder. An HTTP client sends a
query, variables, an operation name and extensions, and has no route to global
state at all, so nothing outside your own process can mark a request this way.
As for the plural in that comment, exactly two middleware in the whole
repository consult the flag: this one, and the one that enforces
persisted-operations-only, which waives the requirement for a warmup request.
One security measure, and it belongs to
chapter~\ref{ch:security-hardening}'s subject rather than this one.

\index{warmup requests!measured}%
Measured on the sample, a warmed executor answers its first client request with
the operation already compiled:

\begin{minted}[fontsize=\footnotesize,breaklines]{text}
== warmup-fills-the-operation-cache
   first request, no warmup task -> document cache hits 0, operation cache hits 0
   first request, warmed         -> document cache hits 0, operation cache hits 1
   PASS
\end{minted}

\index{document cache!needs an identifier to be consulted}%
The document cache stays cold there, and that is not the warmup's doing. It is a
difference between the two harnesses this book has used, and it is worth knowing
before you instrument anything in process. \code{DocumentCacheMiddleware} only
looks in the cache when the request carries a document id or a document hash,
and only stores one when the parsed document ended up with an id. Over HTTP the
transport supplies a hash and this is invisible; a request built in process from
source text carries neither, so there is no lookup to hit. Mosaic's own
timelines report document cache hits all day, which is the same code behaving
differently because the request arrived differently.

\subsection*{The two that arrived with authorization}

\index{authorization!the two middleware, in source}%
One call in \code{AuthorizeRequestExecutorBuilder.cs} registers both, and a
validation rule with them:

\begin{minted}[fontsize=\footnotesize,breaklines]{csharp}
builder.AddValidationRule(
    (s, _) => new AuthorizeValidationRule(
        s.GetRequiredService<AuthorizationCache>()));

var prepareAuthorization = PrepareAuthorizationMiddleware.Create();
builder.UseRequest(
    prepareAuthorization.Middleware,
    key: prepareAuthorization.Key,
    before: WellKnownRequestMiddleware.DocumentValidationMiddleware);

var authorizeRequest = AuthorizeRequestMiddleware.Create();
builder.UseRequest(
    authorizeRequest.Middleware,
    key: authorizeRequest.Key,
    after: WellKnownRequestMiddleware.DocumentValidationMiddleware);
\end{minted}

So chapter~\ref{ch:identity-and-authorization} read the startup log correctly,
and this is the line that put them there, using the named-neighbour mechanism
chapter~\ref{ch:the-life-of-a-request} described for the cost analyzer.

\code{PrepareAuthorizationMiddleware.cs} prepares exactly one thing: it resolves
the authorization handler out of the request scope and stores it in a request
feature. It has to run before validation because the validation rule registered
alongside it runs during validation, and because everything downstream asks for
that handler and throws if nobody put it there.

\index{authorization!refusing by not calling next}%
\code{AuthorizeRequestMiddleware.cs} is the enforcing half, and it refuses by
declining to continue:

\begin{minted}[fontsize=\footnotesize,breaklines]{csharp}
if (result is not AuthorizeResult.Allowed)
{
    context.Result = CreateErrorResult(result);
    return;
}
\end{minted}

That \code{return} is the whole mechanism behind an \code{[Authorize]} in
validation mode refusing a request before any resolver runs: the pipeline stops
there, above the operation cache, the compiler and the executor. The apply-policy
enum has three members. \code{BeforeResolver}, the
default, and \code{AfterResolver} are both a field middleware inserted per
guarded field, checking before or after the resolver runs. \code{Validation}
gets no field middleware at all, and is handled here. Mosaic uses the default
everywhere, which is why chapter~\ref{ch:identity-and-authorization}'s refusals
came back as field errors with a path and an HTTP 200 rather than as a bare
401.

\subsection*{And what that cost, which nobody was watching}

\index{validation!the non-cacheable rule}%
Chapter~\ref{ch:the-life-of-a-request} recorded, as the first of its three
unverified items, that it did not know which shipped validation rules are
non-cacheable, and warned against generalising. Good instinct, wrong worry. The
rule that matters did not exist in this graph then, and
chapter~\ref{ch:identity-and-authorization} installed it.

\begin{minted}[fontsize=\footnotesize,breaklines]{csharp}
public bool IsCacheable => false;
\end{minted}

That is \code{AuthorizeValidationRule.cs}, and the property is hard-coded. It is
not conditioned on any field using \code{apply: Validation}, or on any field
being guarded, or on anything else. It is registered unconditionally by the call
above. And \code{DocumentValidator.cs} computes its verdict once, when it is
built:

\begin{minted}[fontsize=\footnotesize,breaklines]{csharp}
_nonCacheableRules = [.. rules.Where(rule => !rule.IsCacheable)];
...
public bool HasNonCacheableRules => _nonCacheableRules.Length > 0;
\end{minted}

which \code{DocumentValidationMiddleware.cs} then consults on every request:

\begin{minted}[fontsize=\footnotesize,breaklines]{csharp}
if (!documentInfo.IsValidated || _documentValidator.HasNonCacheableRules)
{
    using (_diagnosticEvents.ValidateDocument(context))
    {
        var result =
            _documentValidator.Validate(
                context.Schema,
                documentInfo.Id,
                documentInfo.Document,
                context.Features,
                documentInfo.IsValidated);
\end{minted}

\index{document cache!validation is narrowed, not skipped}%
Two things follow. The smaller one is that
chapter~\ref{ch:the-life-of-a-request} described a binary where the code has a
narrowing: that last argument is \code{onlyNonCacheable}, and on a cache hit it
is true, so a validator with non-cacheable rules re-runs only those rather than
all of them. The larger one is that merely calling \code{AddAuthorization} is
enough to re-open the validation phase on every document-cache hit that service
will ever serve.

\index{validation!measured across the seven subgraphs}%
Which is a claim about the real graph, so here is the real graph. The same
document sent to each of the seven subgraphs three times, and the last timeline
line each of them logged:

\begin{minted}[fontsize=\scriptsize,breaklines,breakanywhere]{text}
catalog     parse - validate - compile - coerce - execute 0.108ms total 0.209ms (document cache hit, operation cache hit, 1 resolvers, 0 SQL)
pricing     parse - validate - compile - coerce - execute 0.101ms total 0.190ms (document cache hit, operation cache hit, 1 resolvers, 0 SQL)
inventory   parse - validate - compile - coerce - execute 0.105ms total 0.205ms (document cache hit, operation cache hit, 1 resolvers, 0 SQL)
accounts    parse - validate 0.037ms compile - coerce - execute 0.090ms total 0.231ms (document cache hit, operation cache hit, 1 resolvers, 0 SQL)
reviews     parse - validate 0.054ms compile - coerce - execute 0.140ms total 0.272ms (document cache hit, operation cache hit, 1 resolvers, 0 SQL)
ordering    parse - validate 0.040ms compile - coerce - execute 0.099ms total 0.238ms (document cache hit, operation cache hit, 1 resolvers, 0 SQL)
nodes       parse - validate - compile - coerce - execute 0.089ms total 0.182ms (document cache hit, operation cache hit, 1 resolvers, 0 SQL)
\end{minted}

Four dashes and three durations. The three are Accounts, Reviews and Ordering,
which are exactly the three that call \code{AddMosaicAuthorization}, and exactly
the three that chapter~\ref{ch:identity-and-authorization} found assembling
fifteen middleware instead of thirteen. One line bought both, and only one of
them was intended.

The claim is the dash rather than the milliseconds. Those durations are
single-machine numbers and move between runs; what is the same everywhere is
whether the phase opened at all, and that is what both verification scripts now
assert.

\subsection*{The run that could have proved me wrong}

\index{authorization!the control}%
That is a correlation across seven services that differ in more than one way, so
it is not yet an argument. The sentence blames \code{AddAuthorization}, so the
run worth making is the one where that call is the only thing that changed.

The obvious control is to take the call out of Accounts, and it is the wrong
one: removing it would leave the \code{[Authorize]} attributes behind and vary
two things. So I went the other way and added
\code{AddMosaicAuthorization()} to Catalog, which has no guarded field, no
attribute and no security registration of any kind, and rebuilt that one
container:

\begin{minted}[fontsize=\scriptsize,breaklines,breakanywhere]{text}
catalog     parse - validate 0.047ms compile - coerce - execute 0.132ms total 0.281ms (document cache hit, operation cache hit, 1 resolvers, 0 SQL)
\end{minted}

Catalog's dash becomes a duration, with nothing in that service to authorise.
The bare registration is sufficient. I put it back immediately, and the same
contrast is committed as a case in \code{samples/executor-internals}, two
executors over one schema differing in that call and nothing else:

\begin{minted}[fontsize=\footnotesize,breaklines]{text}
== authorization-defeats-the-validation-cache
   without AddAuthorization -> 1 validation(s) over 2 requests, 1 document cache hit(s)
   with AddAuthorization    -> 2 validation(s) over 2 requests, 1 document cache hit(s)
   PASS
\end{minted}

Both hit the document cache; only the authorized one validates twice.

\index{validation!the second trigger, outside development}%
There is a second way to acquire a non-cacheable rule and it has nothing to do
with authorization. \code{AddGraphQLServer} disables introspection unless the
host environment is Development, and the rule it uses to do that is registered
as non-cacheable. So any HotChocolate server picks one up the day it runs in
production. Every Mosaic service runs as Development, in the compose file, in
its launch settings and in both verification scripts, so that trigger never
fires here. It would fire in a deployment, and the four services that look
exempt in the capture above are exempt only on this machine.

Finally, the thing I do not want misread. Nothing
chapter~\ref{ch:the-life-of-a-request} wrote was wrong. It measured a graph with
no authorization anywhere and found that a document cache hit skips validation,
and that is still true today of Catalog, Pricing, Inventory and Nodes. What
changed is the graph, in a chapter that was thinking about tokens and scopes.
This is the same shape as the resolver count in
section~\ref{sec:ch16-pure-or-a-task}: a number that stayed correct while the
sentence around it quietly stopped covering every case.

\medskip
Which is the useful thing to take out of a chapter spent inside somebody else's
executor. Not that any of this is wrong, because almost none of it is, but that
the layer you are not reading has opinions about the layer you are, and they
arrive attached to a line you wrote for another reason.
Chapter~\ref{ch:inside-the-router-and-the-composer} does the same walk on the
two processes in front of these seven.

\index{request pipeline!middleware that refuse|)}%
